% Options for packages loaded elsewhere
\PassOptionsToPackage{unicode}{hyperref}
\PassOptionsToPackage{hyphens}{url}
\documentclass[
]{article}
\usepackage{xcolor}
\usepackage[margin=1in]{geometry}
\usepackage{amsmath,amssymb}
\setcounter{secnumdepth}{-\maxdimen} % remove section numbering
\usepackage{iftex}
\ifPDFTeX
  \usepackage[T1]{fontenc}
  \usepackage[utf8]{inputenc}
  \usepackage{textcomp} % provide euro and other symbols
\else % if luatex or xetex
  \usepackage{unicode-math} % this also loads fontspec
  \defaultfontfeatures{Scale=MatchLowercase}
  \defaultfontfeatures[\rmfamily]{Ligatures=TeX,Scale=1}
\fi
\usepackage{lmodern}
\ifPDFTeX\else
  % xetex/luatex font selection
\fi
% Use upquote if available, for straight quotes in verbatim environments
\IfFileExists{upquote.sty}{\usepackage{upquote}}{}
\IfFileExists{microtype.sty}{% use microtype if available
  \usepackage[]{microtype}
  \UseMicrotypeSet[protrusion]{basicmath} % disable protrusion for tt fonts
}{}
\makeatletter
\@ifundefined{KOMAClassName}{% if non-KOMA class
  \IfFileExists{parskip.sty}{%
    \usepackage{parskip}
  }{% else
    \setlength{\parindent}{0pt}
    \setlength{\parskip}{6pt plus 2pt minus 1pt}}
}{% if KOMA class
  \KOMAoptions{parskip=half}}
\makeatother
\usepackage{graphicx}
\makeatletter
\newsavebox\pandoc@box
\newcommand*\pandocbounded[1]{% scales image to fit in text height/width
  \sbox\pandoc@box{#1}%
  \Gscale@div\@tempa{\textheight}{\dimexpr\ht\pandoc@box+\dp\pandoc@box\relax}%
  \Gscale@div\@tempb{\linewidth}{\wd\pandoc@box}%
  \ifdim\@tempb\p@<\@tempa\p@\let\@tempa\@tempb\fi% select the smaller of both
  \ifdim\@tempa\p@<\p@\scalebox{\@tempa}{\usebox\pandoc@box}%
  \else\usebox{\pandoc@box}%
  \fi%
}
% Set default figure placement to htbp
\def\fps@figure{htbp}
\makeatother
\setlength{\emergencystretch}{3em} % prevent overfull lines
\providecommand{\tightlist}{%
  \setlength{\itemsep}{0pt}\setlength{\parskip}{0pt}}
\usepackage{bookmark}
\IfFileExists{xurl.sty}{\usepackage{xurl}}{} % add URL line breaks if available
\urlstyle{same}
\hypersetup{
  hidelinks,
  pdfcreator={LaTeX via pandoc}}

\author{}
\date{\vspace{-2.5em}}

\begin{document}

\section{Govorne bilješke za
prezentaciju}\label{govorne-biljeux161ke-za-prezentaciju}

Trajanje: ciljaj 12-14 minuta i ostavi 1-2 minute za pitanja. Govori
mirno; gotovo svaka slajd ima jednu glavnu poruku.

\subsection{Slajd 1 --- Naslov}\label{slajd-1-naslov}

Reci: Tema je evaluacija jezičnih generativnih modela za
audiorehabilitaciju. Cilj nije samo generirati lijepe riječi, nego
provjeriti mogu li se dobiti hrvatske riječi i kratke rečenice s
kontroliranim fonemskim sastavom, a zatim ih pretvoriti u audio.

\subsection{Slajd 2 --- Motivacija}\label{slajd-2-motivacija}

Objasni problem: u audiorehabilitaciji trebamo materijal koji cilja
određene glasove ili kontraste. Ako korisnik ima problem razlikovati /s/
i /š/, nije dovoljno imati bilo koju rečenicu. Treba znati koji fonemi
se stvarno pojavljuju i u kojem omjeru.

\subsection{Slajd 3 --- Istraživačko
pitanje}\label{slajd-3-istraux17eivaux10dko-pitanje}

Glavna ideja: LLM može biti koristan kao generator kandidata, ali nije
pouzdan evaluator. Zato je doprinos rada reproducibilni Python pipeline
koji sve broji i provjerava deterministički.

Ako profesor pita „zašto ne vjerovati ChatGPT-u?'': Zato što LLM često
daje uvjerljiv odgovor, ali nema garantirano točno brojanje fonema,
pogotovo za hrvatske višeslovne foneme dž, lj i nj.

\subsection{Slajd 4 --- Veza s
literaturom}\label{slajd-4-veza-s-literaturom}

Reci da se oslanjamo na rad Andrijašević i Vukelić, gdje se GPT koristi
za generiranje hrvatskog govornog materijala za auditivni trening. Mi
preuzimamo ideju pet klasa fonema i razina zasićenja, ali uvodimo strožu
validaciju.

Važno: HJP u promptu je samo uputa modelu. U našem radu HJP valjanost je
zasebna ručna provjera riječi.

\subsection{Slajd 5 --- Fonemske klase i
formula}\label{slajd-5-fonemske-klase-i-formula}

Objasni formulu: broj fonema ciljne klase dijeli se s ukupnim brojem
fonema i množi sa 100. Ako ciljamo klasu N na 70 \%, kandidat prolazi
samo ako barem 70 \% njegovih fonema pripada klasi N.

Tehnički detalj: parser ide longest-match-first za dž, lj i nj. Bez toga
bi npr. „panj'' pogrešno postao p-a-n-j umjesto p-a-nj.

\subsection{Slajd 6 --- Pipeline}\label{slajd-6-pipeline}

Objasni što skripta radi:

\begin{enumerate}
\def\labelenumi{\arabic{enumi}.}
\tightlist
\item
  Učita CSV ili generira kandidate preko Ollame.
\item
  Normalizira tekst u UTF-8, čuva hrvatske znakove.
\item
  Fonemizira tekst.
\item
  Računa zasićenje.
\item
  Provjerava dopuštene znakove, broj riječi, duplikate i ponovljene
  riječi.
\item
  Dodaje Hunspell screening i ručni HJP word-review.
\item
  Stvara CSV i Markdown izvještaje.
\end{enumerate}

Ako pita „koje su glavne datoteke'': \texttt{src/phonemizer.py},
\texttt{src/validators.py}, \texttt{src/metrics.py},
\texttt{src/pipeline.py}, \texttt{src/tts.py},
\texttt{src/asr\_eval.py}.

\subsection{Slajd 7 --- Dizajn
eksperimenata}\label{slajd-7-dizajn-eksperimenata}

Reci da postoje tri sloja:

Prvi je reprodukcija paper-style promptova. Drugi je Task 16 usporedba
ChatGPT Plus protiv lokalne Ollame. Treći je audio faza samo na
validiranim kandidatima.

Naglasak: audio ne radimo nad svim kandidatima, nego samo nad onima koji
su prošli determinističku i leksičku provjeru.

\subsection{Slajd 8 --- Reprodukcija rada:
riječi}\label{slajd-8-reprodukcija-rada-rijeux10di}

Ovdje reci da je ovo namjerno prikazano gotovo istim formatom kao Table
II u radu. U paperu su retci fonemske klase, stupci su saturation
leveli, a ćelija je postotak riječi koje zadovoljavaju oba kriterija:
zasićenje + standardni hrvatski/HJP.

Naša tablica koristi isti princip: saturation pass + ručni HJP
word-review. Važno je reći da naš pipeline dodatno mjeri duplikate, ali
ih ovdje ne miješamo u ćeliju jer želimo metodološki isti indikator kao
u paperu.

Zaključak za riječi: naši paper-style rezultati su vrlo visoki za većinu
klasa, ali SV na 70 \% i V na 80 \% pokazuju pad. To je usporedivo s
idejom rada da visoka zasićenja i ``visoke'' klase postaju teže.

\subsection{Slajd 9 --- Reprodukcija rada:
rečenice}\label{slajd-9-reprodukcija-rada-reux10denice}

Ovdje je analog Table III iz rada. Paper za rečenice prikazuje postotak
rečenica koje zadovoljavaju saturation level. Mi prikazujemo isti
kriterij.

Glavna poruka: na 50 \% smo gotovo savršeni za većinu klasa. Na 70 \% se
vidi pad, posebno za klasu V. To je važan zaključak jer se slaže s
očekivanjem: viši prag daje modelu manje slobode.

\subsection{Slajd 10 --- ChatGPT Plus vs
Ollama}\label{slajd-10-chatgpt-plus-vs-ollama}

Ovo je najvažniji tekstni rezultat. ChatGPT Plus ima 48,6 \% tehnički
valjanih kandidata i 85,9 \% saturation pass. Ollama ima samo 4,5 \%
tehnički valjanih i 4,6 \% saturation pass.

Zaključak: u ovoj konfiguraciji lokalni \texttt{llama3.1:8b} nije
dovoljan za strogo hrvatsko fonemsko zasićenje. ChatGPT Plus je znatno
bolji generator, ali i dalje treba validaciju.

\subsection{Slajd 11 --- Analiza
grešaka}\label{slajd-11-analiza-greux161aka}

Objasni razliku u vrstama grešaka:

ChatGPT Plus: najveći problem su duplikati. To znači da model često
razumije zadatak, ali se vrti oko istih sigurnih kandidata.

Ollama: najveći problem je sam fonetski kriterij. To je ozbiljniji
problem jer znači da generirani tekst ne zadovoljava glavnu znanstvenu
mjeru.

\subsection{Slajd 12 --- Tehnička valjanost i
PCD}\label{slajd-12-tehniux10dka-valjanost-i-pcd}

Ovo je slajd za metodološko objašnjenje. Reci da \texttt{is\_valid} nije
procjena LLM-a, nego rezultat determinističkih pravila.

Kandidat je tehnički valjan ako prođe saturation threshold, nema
nedopuštene znakove, ima ispravan broj riječi, nije duplikat i nema
ponovljene riječi u rečenici. U glavnoj usporedbi Hunspell i HJP držimo
odvojeno jer su leksička, a ne čisto fonetsko-strukturna provjera.

PCD znači Phonetic Content Dissimilarity. Naš paper-style PCD uspoređuje
fonemski sadržaj dvaju kandidata: gleda koliko fonema u duljem kandidatu
nije dijeljeno s kraćim, podijeljeno duljinom duljeg kandidata. Veći
prosječni PCD znači raznolikiji skup.

Rezultat: ChatGPT Plus ima viši prosječni PCD po grupama od Ollame, što
znači da je osim više validnih kandidata dao i korisniju fonetsku
raznolikost. Ollama ima malo grupa s dovoljno validnih kandidata, pa je
i PCD slabiji.

\subsection{Slajd 13 --- Hunspell i HJP}\label{slajd-13-hunspell-i-hjp}

Reci da Hunspell nije HJP. Hunspell je automatski screening: može
odbaciti valjane flektirane riječi ili prihvatiti riječ koja nije
prikladna za rehabilitacijski kontekst.

Ručni HJP word-review: iz kandidata se izvezu jedinstvene riječi, označi
se \texttt{hjp\_valid}, a zatim se ta odluka propagira natrag na riječi
i rečenice. Rečenica je HJP-valid samo ako su sve njezine riječi
HJP-valid.

\subsection{Slajd 14 --- Audio dizajn}\label{slajd-14-audio-dizajn}

Objasni da je za poštenu TTS usporedbu isti skup od 96 tekstova poslan u
sva tri sustava: eSpeak NG, Coqui VITS HR i SpeechT5 HR. Time se ne
uspoređuju različiti tekstovi, nego različiti sintetizatori na istom
materijalu.

Svi izlazi su normalizirani u isti format: WAV, mono, 16 kHz, 16-bit
PCM.

\subsection{Slajd 15 --- TTS i ASR}\label{slajd-15-tts-i-asr}

Sva tri TTS sustava tehnički su uspjela sintetizirati 96/96 kandidata.
ASR evaluacija koristi fiksni \texttt{faster-whisper\ large-v3-turbo}
profil za hrvatski.

Reci caveat: WER/CER nije klinički dokaz. To je automatski proxy. Visok
WER može značiti da je TTS loš, ASR loš, ili da je materijal fonemski
neobičan.

\subsection{Slajd 16 --- Audio demo}\label{slajd-16-audio-demo}

Ovdje pusti tri WAV datoteke istim redom: eSpeak NG, Coqui VITS HR,
SpeechT5 HR. Naglasi da je tekst isti: „Draga rada radi.'' Time se sluša
razlika u sintetizatoru, a ne razlika u sadržaju.

Reci da su demo kopije pojačane za prezentaciju jer je SpeechT5 u
originalnom izlazu bio tiši. Originalni istraživački audio nije
mijenjan; pojačanje je samo praktično za slušanje u učionici.

\subsection{Slajd 17 --- Slušna
provjera}\label{slajd-17-sluux161na-provjera}

Ovo je važan ljudski rezultat. eSpeak NG je najrazumljiviji po ljudskoj
ocjeni, iako nije najprirodniji. SpeechT5 je najprirodniji, ali ne
najrazumljiviji. Coqui je u ovom testu najslabiji.

Poanta: automatska ASR metrika i ljudska procjena ne moraju dati isti
poredak, zato se ne smije stati samo na WER/CER.

\subsection{Slajd 18 --- Zaključak}\label{slajd-18-zakljuux10dak}

Zaključi:

\begin{enumerate}
\def\labelenumi{\arabic{enumi}.}
\tightlist
\item
  Pipeline radi i reproducibilan je.
\item
  LLM može pomoći u generiranju kandidata, ali samo ako Python
  deterministički validira.
\item
  ChatGPT Plus je bio bolji od lokalne Ollame za tekst.
\item
  Audio faza je tehnički izvediva, ali klinička prikladnost traži
  stručnu provjeru.
\end{enumerate}

\subsection{Kratki odgovori ako profesor
pita}\label{kratki-odgovori-ako-profesor-pita}

\textbf{Kako fonemizator radi?}\\
Tekst se pretvara u mala slova, uklanja se interpunkcija, čuvaju se
hrvatski znakovi, a \texttt{dž}, \texttt{lj} i \texttt{nj} se prepoznaju
prije pojedinačnih slova. Razmaci ne ulaze u broj fonema.

\textbf{Što znači saturation pass?}\\
Ako kandidat ima N fonema, a nk pripada ciljnoj klasi, zasićenje je nk/N
× 100. Kandidat prolazi ako je postotak veći ili jednak traženom pragu.

\textbf{Zašto imate Hunspell i HJP?}\\
Hunspell je automatski i skalabilan, ali nije konačan. HJP/manual
word-review je bliže kriteriju iz literature, ali je ručni.

\textbf{Zašto je Ollama tako slaba?}\\
Testirani lokalni model nije pouzdano pratio hrvatske fonemske klase.
Većina kandidata pala je na zasićenju, ne samo na hrvatskoj leksici.

\textbf{Zašto TTS nije klinički zaključak?}\\
Jer tehnički WAV i dobra ASR transkripcija ne dokazuju da je izgovor
terapeutski prikladan. Potrebni su stručnjaci, više slušatelja i
eventualno korisničko testiranje.

\textbf{Što biste poboljšali?}\\
Veći ručni listening review, više ocjenjivača, inter-rater reliability,
jači lokalni LLM, repair loop za neuspjele kandidate i formalniji
HJP/klinički review.

\end{document}
